\chapter*{How to use this booklet}
\addcontentsline{toc}{chapter}{How to use this booklet}
\markboth{How to use this booklet}{How to use this booklet}

This booklet accompanies the hands-on tutorial on the new SHiP software
framework. It is built as a sequence of small steps (\cref{fig:roadmap}): every
chapter says at its start what it needs and what you will be able to do
afterwards, and ends with a short self-check. Each step uses only what the
steps before it taught.

Part~I is the way of working, with no SHiP software yet: git on your own, then
with branches, then with other people and pull requests; pixi environments;
continuous integration; and finally a small analysis that the whole group
writes together, one reviewed and CI-checked pull request per person.
Part~II follows one physics example from beginning to end, from scratch: the
muons that 400\,GeV protons produce in the SHiP target, the largest
background of the experiment. You generate the proton collisions yourself and
see what leaves the hadron absorber, then follow an artificial sample of
muons from the back of the muon shield through the detector. Every chapter of Part~II adds one piece of the software chain, and
each piece consumes what the previous one produced.

\begin{figure}[ht]
\centering
\begin{tikzpicture}[
  rstep/.style={draw=shipnavy, fill=white, rounded corners=2pt, minimum height=8mm,
    text width=21mm, align=center, font=\sffamily\scriptsize, inner sep=2pt,
    execute at begin node={\hyphenpenalty=10000\exhyphenpenalty=10000}},
  capstone/.style={rstep, fill=shipcream, line width=0.9pt},
  next/.style={-{Stealth[length=1.6mm]}, color=shipnavy},
  partlabel/.style={font=\sffamily\small\bfseries, anchor=west, color=shipnavy}]
  % Part I: two rows of three
  \node[partlabel] at (-1.25,0.85) {Part I: working together};
  \node[rstep] (g1) at (0,0)   {\ref{ch:git}\enspace git on your own};
  \node[rstep] (g2) at (2.6,0) {\ref{ch:branches}\enspace branches, merge, rebase};
  \node[rstep] (g3) at (5.2,0) {\ref{ch:remotes}\enspace remotes, pull requests};
  \node[rstep] (px) at (7.8,0) {\ref{ch:pixi}\enspace pixi environments};
  \node[rstep] (ci) at (10.4,0) {\ref{ch:ci}\enspace continuous integration};
  \node[capstone] (tm) at (13.0,0) {\ref{ch:team}\enspace team analysis};
  \foreach \a/\b in {g1/g2, g2/g3, g3/px, px/ci, ci/tm} \draw[next] (\a) -- (\b);
  % Part II: two rows
  \node[partlabel] at (0.35,-1.15) {Part II: the SHiP software chain};
  \node[rstep] (s1) at (0,-2)    {\ref{ch:setup}\enspace SHiP environment};
  \node[rstep] (s2) at (2.6,-2)  {\ref{ch:bigpicture}\enspace the big picture};
  \node[rstep] (s3) at (5.2,-2)  {\ref{ch:geometry}\enspace Geometry};
  \node[rstep] (s4) at (7.8,-2)  {\ref{ch:services}\enspace geometry and field services};
  \node[rstep] (s5) at (10.4,-2) {\ref{ch:datamodel}\enspace data model};
  \node[rstep] (s6) at (13.0,-2) {\ref{ch:production}\enspace fixed-target production};
  \node[rstep] (s7) at (13.0,-3.3) {\ref{ch:aegir}\enspace muons in the detector};
  \node[rstep] (s8) at (10.4,-3.3) {\ref{ch:display}\enspace \mbox{sea-cucumber} display};
  \node[rstep] (s9) at (7.8,-3.3)  {\ref{ch:digitisation}\enspace Shannon digitisation};
  \node[capstone] (s10) at (5.2,-3.3) {\ref{ch:analysis}\enspace muon analysis};
  \node[rstep, dashed] (s11) at (2.6,-3.3) {\ref{ch:outlook}\enspace Trout: discussion};
  \foreach \a/\b in {s1/s2, s2/s3, s3/s4, s4/s5, s5/s6, s6/s7, s7/s8, s8/s9, s9/s10, s10/s11}
    \draw[next] (\a) -- (\b);
  \draw[next, dashed] (tm.south) -- ++(0,-0.35) -| (s1.north);
\end{tikzpicture}
\caption{The steps of the tutorial, numbered by chapter. Each step builds on
the one before it; the shaded steps are the two places where everything
learnt so far comes together. Part~II starts once Part~I is done.}
\label{fig:roadmap}
\end{figure}

\begin{pitfall}
This is a reviewed teaching draft, not an end-to-end validated release.
The archive contains no field maps and no frozen common environment, and the
supplied workflows have not been run end to end. The small samples of the
exercises illustrate the interfaces; their rates are not a background
prediction. See \cref{app:review} for the checked source revisions and
remaining prerequisites.
\end{pitfall}

\section*{What you will learn}

\begin{itemize}
  \item how to work with git and GitHub the way the ShipSoft repositories
        do: commits, branches, merging versus rebasing, remotes, forks and
        pull requests with review (\cref{ch:git,ch:branches,ch:remotes});
  \item how to describe and reproduce a software environment with pixi:
        manifest, lock file, tasks and environments (\cref{ch:pixi});
  \item how continuous integration checks every pull request, with GitHub
        Actions as in the ShipSoft repositories (\cref{ch:ci});
  \item how to build an analysis as a team, one reviewed, CI-checked pull
        request per person, by fitting a mystery resonance with PyROOT
        (\cref{ch:team});
  \item how the SHiP detector is described in \pkg{GeoModel}
        (\repo{Geometry}) and how other packages reach it through the
        \pkg{geometry service} (\repo{geometry\_service});
  \item how magnetic field maps are stored and served
        (\repo{field\_service});
  \item what the event data model looks like on disk (\repo{data-model});
  \item how to simulate proton collisions in the target and the muons they
        produce with \pkg{aegir} (\repo{aegir}), the Phlex-based simulation
        built on Pythia\,8 and Geant4: a production that ends on a scoring
        plane behind the hadron absorber, and muons through the detector;
  \item how to look at them with the event display \pkg{sea-cucumber}
        (\repo{sea-cucumber});
  \item how to digitise them with \pkg{Shannon} (\repo{Shannon});
  \item how to analyse the resulting files directly from Python.
\end{itemize}

Reconstruction (\pkg{Trout}) is not part of the hands-on programme; it is the
subject of the closing discussion (\cref{ch:outlook}).

\section*{Conventions}

Commands to type in a terminal look like this:
\begin{shell}
pixi cvmfs phlex -c <(jsonnet --tla-code events=100 ft_production.jsonnet)
\end{shell}
Illustrative output is shown with a vertical rule on the left; exact log text may differ:
\begin{outputlisting}
[info] sim_output_module: writing 'ft_production.root' ...
\end{outputlisting}

Package names are set in \pkg{bold sans serif}, files as \file{path/to/file},
and data-model fields as \field{pdg\_code}. Lengths are in millimetres,
momenta and energies in GeV, and times in nanoseconds throughout, matching the
software. The coordinate system is the one of the SHiP geometry: the origin is
the front face of the target, $z$ points downstream along the beam, $y$
points up (\cref{sec:coordinates}).

The boxes have fixed meanings:

\begin{exercise}{Something for you to do}
Exercises are numbered by chapter. Most take five to fifteen minutes.
\end{exercise}

\begin{checkpoint}
What you should see when the exercise worked. If you do not, ask a tutor
before moving on: later chapters use the files you make here.
\end{checkpoint}

\begin{note}
Background, or where to read more.
\end{note}

\begin{pitfall}
A mistake that is easy to make, and how to recognise it.
\end{pitfall}

\iforganiser
\begin{organisernote}
This is the organiser edition. It contains the solutions to every exercise
and notes for the tutors, in boxes like this one. The participant edition is
built from the same sources without them (\file{booklet.tex}; this edition is
\file{booklet-organiser.tex}). Items still to be decided are marked
\tbd{like this} in both editions; \cmd{make check} lists them.
\end{organisernote}
\fi

\section*{The files you need}

The runnable configurations and analysis scripts are supplied in the
\file{framework-booklet/code/} directory of the accompanying archive.
Inline API excerpts are explanatory. Extract the archive, then run:

\begin{shell}
cd framework-booklet
export CODE="$PWD/code"
export SHIP_TUTORIAL="$(dirname "$PWD")/ship-tutorial"
mkdir -p "$SHIP_TUTORIAL/work"
\end{shell}
The commands in this booklet refer to that directory as \texttt{\$CODE}; set
both variables again in every new terminal. Use Bash for commands containing
\texttt{<(...)}, and exit a package shell before activating another.

\section*{Getting help}

Each package documents its classes in a Doxygen reference linked from its
README, and the collaboration-wide documentation lives at
\url{https://shipsoft.github.io/Documentation/}. Please report mistakes in this
booklet to \href{mailto:matclim@cern.ch}{matclim@cern.ch}.
